\clearpage
\VerdierPhysicalPageAnchor{0196}%
\begingroup\small\setlength{\parskip}{0pt}\setlength{\emergencystretch}{1.5em}
\setlength{\abovedisplayskip}{4pt plus 2pt minus 2pt}%
\setlength{\belowdisplayskip}{4pt plus 2pt minus 2pt}%

\VerdierSemanticAnchor{ch3d-2-3-1-construction}
Ceci étant fait, il suffira de poser
$L^\bullet=\varprojlim_n L_n^\bullet$,
$w=\varprojlim_n w_n$. Pour $n$ suffisamment petit, on pose
$L_n=Y$, $v_n=\operatorname{id}_Y$ et
$w_n=\operatorname{id}_Y$. Supposons que l'on ait construit pour tout entier
$n<q$, $q\in\mathbf Z$, des complexes $L_n^\bullet$ et des morphismes $v_n$
et $w_n$ possédant les propriétés 1), 2) et 3) ci-dessus. Le complexe $X_n$
est de la forme :

\VerdierSemanticAnchor{ch3d-2-3-1-xn-shape}
\[
X_n=\cdots\longrightarrow Y^{-n-2}\longrightarrow Z^{-n-1}
\longrightarrow X^{-n}\longrightarrow X^{-n+1}\longrightarrow\cdots
\]
et le complexe $L_n^\bullet$ est de la forme :

\VerdierSemanticAnchor{ch3d-2-3-1-ln-shape}
\[
L_n^\bullet=\cdots\longrightarrow Y^{-n-2}\longrightarrow Z'^{-n-1}
\longrightarrow L^{-n}\longrightarrow L^{-n+1}\longrightarrow\cdots,
\]
où les $Y^i$ (resp.\ les $X^i$) sont les composants du complexe $Y$
(resp. $X$). On a donc un diagramme :

\VerdierSemanticAnchor{ch3d-2-3-1-degree-q-diagram}
\begin{center}
\begin{tikzcd}[column sep=5.2em,row sep=2.8em]
  & X^{-q} \arrow[d,"(\dot u_q)^{-q}"] \\
Z'^{-q} \arrow[r,"(w_{q-1})^{-q}"] & Z^{-q},
\end{tikzcd}
\end{center}
où $(\dot u_q)^{-q}$ est un épimorphisme et où $Z'^{-q}\in M^\wedge$.
Comme l'ensemble $M^\wedge$ possède la propriété $(E_2)$ (2.2.4), on peut
compléter ce diagramme en un diagramme commutatif :

\VerdierSemanticAnchor{ch3d-2-3-1-lifting-square}
\begin{center}
\begin{tikzcd}[column sep=5.2em,row sep=2.8em]
L^{-q} \arrow[r,"\beta"] \arrow[d,"\alpha"']
  & X^{-q} \arrow[d,"(\dot u_q)^{-q}"] \\
Z'^{-q} \arrow[r,"(w_{q-1})^{-q}"] & Z^{-q},
\end{tikzcd}
\end{center}
où le morphisme $\alpha$ est un épimorphisme et l'objet $L^{-q}$ un élément
de $M$. On en déduit, en utilisant les notations (2.1.1.2) et (2.1.1.3), un
diagramme commutatif de complexes :

\VerdierSemanticAnchor{ch3d-2-3-1-complex-square}
\begin{center}
\begin{tikzcd}[column sep=4.8em,row sep=3.0em]
L_{q-1}^\bullet(\alpha,-q) \arrow[r,"\varphi"]
  \arrow[d,"{p(\alpha,-q)}"']
  & X_{q-1}((\dot u_q)^{-q},-q)\simeq X_q
    \arrow[d,"{p((\dot u_q)^{-q},-q)\simeq\dot u_q}"] \\
L_{q-1}^\bullet \arrow[r,"w_{q-1}"] & X_{q-1}.
\end{tikzcd}
\end{center}
\endgroup
